\documentclass[11pt,a4paper]{article}

\usepackage{geometry}
\geometry{margin=2.5cm}

\usepackage[T1]{fontenc}
\usepackage{amsmath,amssymb,amsthm}
\usepackage{hyperref}
\usepackage{booktabs}
\usepackage{graphicx}
\usepackage{xcolor}
\usepackage{titlesec}
\usepackage{enumitem}

\hypersetup{
  colorlinks=true,
  linkcolor=blue!70!black,
  citecolor=green!50!black,
  urlcolor=blue!60!black
}

\titleformat{\section}{\Large\bfseries}{\thesection}{1em}{}
\titleformat{\subsection}{\large\bfseries}{\thesubsection}{1em}{}

\newtheorem{remark}{Remark}

\title{\textbf{EAS vs.\ Friston Free Energy Principle:\\A Deep Comparative Analysis}}
\author{Jing Zhang\\[4pt]
\small Independent Researcher\\
\small ORCID: 0009-0008-3136-2457}
\date{July 2, 2026}

\begin{document}
\maketitle

\begin{abstract}
The Epistemic Axiomatic System (EAS) and Karl Friston's Free Energy Principle (FEP) represent two ambitious frameworks for understanding intelligence and cognition. This paper provides a systematic comparison, establishing that the two frameworks are ``orthogonal but in dialogue'': FEP describes how cognitive systems operate (minimize free energy), while EAS delineates what cognitive systems cannot achieve (perfect self-modeling is structurally impossible). We identify five core divergences, trace the latest Friston literature (2024--2026), and propose that EAS's single axiom can provide first-principles grounding for FEP's claim that approximate inference is necessarily approximate.
\end{abstract}

\tableofcontents
\newpage

\section{Summary}

EAS's single axiom $S \subset E \implies R_S \not\cong E$ (finite observation is necessarily a quotient of reality; perfect self-modeling is impossible) and Friston's variational free energy minimization stand in a relationship that is \textbf{neither purely complementary nor purely competitive, but ``orthogonal but in dialogue.''} FEP proceeds from dynamics, describing how cognitive systems operate; EAS proceeds from epistemology, delineating what cognitive systems cannot achieve.

The key divergence: FEP frames cognitive limitations as an \textbf{engineering problem} (approximate inference can be progressively improved), while EAS frames them as a \textbf{structural theorem} (perfect self-modeling is impossible even with unlimited computation).

As of mid-2026, Friston has not formally theorem-ized G\"{o}delian incompleteness or ``perfect self-modeling impossibility'' in published work---although the 2025 Fields--Friston paper suggested that systems cannot have fully precise models of the QRFs they interact with (paraphrase, not a verbatim quote) \cite{fields2025}. The optimal entry point for dialogue is: \textbf{using EAS's single axiom to provide a first-principles proof of why FEP's approximate inference can only be approximate}.

\section{Does FEP Implicitly Contain EAS's Quotient Cognition Constraint?}

\subsection{FEP's Core Assumptions}

FEP's argument chain can be summarized in three steps:

\begin{enumerate}[leftmargin=*]
\item \textbf{Markov blanket demarcation}: Any system that persists over time necessarily possesses a Markov blanket, rendering internal states conditionally independent of external states \cite{friston2019}.

\item \textbf{Variational free energy upper bound}: Internal states can be interpreted as approximate encodings of posterior distributions; variational free energy is a computable upper bound on surprisal \cite{friston2013}.

\item \textbf{Self-evidencing}: System dynamics necessarily manifest as free energy minimization---not by choice, but as a consequence of existence.
\end{enumerate}

Crucially, FEP explicitly employs \textbf{approximate} Bayesian inference, acknowledging that exact inference is computationally intractable \cite{friston2013}. The 2025 Friston--Sakthivadivel paper further clarifies that FEP is an ``as if'' description---``using this framework need not presuppose that the system is truly performing inference; it is a useful explanatory fiction'' \cite{friston2025}.

\subsection{FEP's Approximation vs.\ EAS's Impossibility}

\begin{center}
\begin{tabular}{@{}lll@{}}
\toprule
\textbf{Dimension} & \textbf{FEP} & \textbf{EAS} \\
\midrule
Nature of cognitive limits & Computational constraint (tractability) & Structural impossibility (theorem) \\
Terminal state of approx.\ inference & Can asymptotically approach true posterior & True posterior is never reachable (quotient constraint) \\
Self-modeling & ``Beautiful Loop'' enables recursive self-reference & Perfect self-modeling is theorem-level impossible \\
Source of limitation & Biological/computational resource limits & Epistemic structure ($S \subset E$ subset relation) \\
\bottomrule
\end{tabular}
\end{center}

\textbf{Core judgment}: FEP implicitly acknowledges cognitive limitations but frames them as \textbf{improvable approximations} rather than \textbf{insurmountable theorems}. FEP's variational free energy $F \geq -\ln p(o)$ allows equality when the variational density equals the true posterior---meaning FEP's mathematical structure \textbf{preserves the ideal observer as a limit}. EAS's single axiom asserts that this limit is \textbf{epistemically unreachable}---not ``hard to reach,'' but ``structurally impossible.''

\section{Has Friston (2024--2026) Addressed ``Incompleteness''?}

\subsection{Recent Key Papers}

\begin{center}
\begin{tabular}{@{}llp{8cm}@{}}
\toprule
\textbf{Paper} & \textbf{Year} & \textbf{Relation to Incompleteness} \\
\midrule
Friston, Ramstead \& Sakthivadivel, ``Generative modelling in non-equilibrium statistical mechanics'' & 2025 & Distinguishes model from modeled system; does not discuss self-modeling impossibility \\
Fields, Albarracin, Friston et al., ``Inner screens and imaginative experience'' & 2025 & \textbf{Closest}: suggests systems cannot have fully precise models of their QRFs (paraphrase), but as practical observation, not theorem \\
Laukkonen, Friston \& Chandaria, ``A Beautiful Loop'' & 2025 & Discusses self-reference and consciousness, but frames as positive emergence (self-reference produces consciousness), not negative constraint (self-reference cannot be perfect) \\
Five Fristonian Formulae (Entropy special issue) & 2025 & Review; does not address incompleteness \\
\bottomrule
\end{tabular}
\end{center}

\subsection{Clear Conclusion}

\textbf{As of mid-2026, Friston has not formally addressed G\"{o}delian incompleteness or cognitive boundary problems.} Fields--Friston (2025) discussed that systems cannot have fully-accurate models of the QRFs they interact with \cite{fields2025}---this can be understood as a consequence of Ashby's Required Variety Law (this is our interpretation, not an explicit attribution in the original text), but it is not elevated to an epistemic axiom or insurmountable theorem. The ``Beautiful Loop'' touches on self-referential structure, but its core thesis is \textbf{how self-reference produces consciousness}, not \textbf{why self-reference cannot be perfectly closed}.

\section{Epiplexity Five-Layer Resonance vs.\ FEP Variational Free Energy}

\subsection{Which Is More Fundamental?}

EAS's Epiplexity five-layer resonance validation (internal $\to$ cross-domain $\to$ predictive $\to$ compressive $\to$ fractal) is a \textbf{validation methodology}---it does not claim to describe how cognition operates, but rather that \textbf{any valid epistemic conclusion must pass five-layer resonance to count as validated}. FEP's variational free energy is a \textbf{descriptive principle}---it claims cognitive systems necessarily minimize this quantity.

The two do not compete at the same level:

\begin{itemize}[leftmargin=*]
\item \textbf{Variational free energy} answers: What do cognitive systems do? (Minimize free energy.)
\item \textbf{Epiplexity five-layer resonance} answers: How do we know an epistemic conclusion is valid? (Must pass five layers; one layer is insufficient.)
\end{itemize}

\subsection{Which Is More Operational?}

\begin{center}
\begin{tabular}{@{}lll@{}}
\toprule
\textbf{Dimension} & \textbf{Epiplexity Five-Layer} & \textbf{Variational Free Energy} \\
\midrule
Mathematical definition & Five-layer validation protocol & Precisely defined: $F = D_{KL}[q\|p] - \ln p(o)$ \\
Computability & Layer-by-layer judgment; no closed formula & Computable via predictive coding etc. \\
Empirical grounding & Each layer requires independent validation & Extensive predictive coding evidence in neuroscience \\
Scope & Validation of epistemic conclusions & Any Markov blanket system \\
\bottomrule
\end{tabular}
\end{center}

\textbf{Judgment}: Variational free energy far surpasses in operationality---it has closed mathematical form, computable approximations, and neuroscience implementation paths. Epiplexity five-layer resonance, as a validation framework, provides value in \textbf{preventing epistemic overreach}, but lacks the computational tractability of variational free energy.

\section{Core Divergences}

\subsection{Divergence 1: Cognitive Limits as ``Engineering Approximation'' vs.\ ``Structural Impossibility''}

\begin{itemize}[leftmargin=*]
\item \textbf{FEP}: Approximate inference exists because exact inference is computationally intractable; in principle, variational density can approach the true posterior arbitrarily closely.
\item \textbf{EAS}: $S \subset E \implies R_S \not\cong E$---even with unlimited computation, perfect self-modeling remains impossible; this is an epistemic-structural necessity.
\item \textbf{Friston's possible response}: FEP's ``as if'' framework already does not commit to the system possessing the true posterior, so EAS's constraint is already implicitly acknowledged within FEP.
\item \textbf{EAS's rebuttal}: FEP's mathematical structure preserves $q = p(\cdot|o)$ as a limit (KL divergence can approach zero), implying the ideal observer is conceivable as an asymptotic limit. EAS explicitly denies this.
\end{itemize}

\subsection{Divergence 2: Imperfect Closure of Self-Reference---EAS Provides Formal Grounding for Beautiful Loop}

\begin{itemize}[leftmargin=*]
\item \textbf{FEP} (``Beautiful Loop''): Self-referential recursion produces consciousness---the world model ``knows'' itself non-locally \cite{laukkonen2025}. \textbf{Importantly}: Beautiful Loop itself does not require perfect closure---its core feature is ``field-evidencing,'' i.e., consciousness arises from the \emph{non-perfectly-closed} state of self-referential loops, not from perfect self-consistency.
\item \textbf{EAS's complement}: EAS does not \emph{contradict} Beautiful Loop but provides the formal foundation that Beautiful Loop lacks. EAS's single axiom $R_S \not\cong E$ rigorously proves that self-referential loops are \textbf{necessarily} imperfect---any self-model is necessarily a quotient space, and residual information loss cannot be eliminated. This provides mathematical grounding for Beautiful Loop's ``non-local knowing'': consciousness arises not \emph{despite} the gap, but \emph{because} the gap is structurally necessary.
\item \textbf{Key difference}: Beautiful Loop describes what this imperfect closure \emph{produces} (consciousness); EAS explains why this imperfect closure \emph{must exist} (quotient cognition constraint). The two are descriptive vs.\ explanatory---not adversarial.
\end{itemize}

\subsection{Divergence 3: Markov Blanket as ``Enabler of Cognition'' vs.\ ``Cage of Cognition''}

\begin{itemize}[leftmargin=*]
\item \textbf{FEP}: Markov blanket makes inference possible---without boundaries, there is no inferring agent.
\item \textbf{EAS}: Markov blanket is precisely the realization of $S \subset E$---internal states can only be approximate encodings of external states; the blanket itself is the boundary of the quotient map.
\item \textbf{Friston's possible response}: This is exactly FEP's insight---cognition is achieved through blanket conditional independence, not limited by it.
\item \textbf{EAS's rebuttal}: FEP sees only the enabling face, not the limiting face. EAS's single axiom proves that enabling and limiting are two sides of the same coin.
\item \textbf{Related critical literature}: Raja et al.\ (2021) ``The Markov Blanket Trick'' and Bruineberg et al.\ (2022) ``The Emperor's New Markov Blankets'' \cite{bruineberg2022} critically examine the metaphysical status of Markov blankets from different angles---the former argues that blanket demarcation may be an epistemic convenience rather than an ontological fact; the latter distinguishes ``Pearl blankets'' (epistemic tools in Bayesian networks) from ``Friston blankets'' (physical boundaries in FEP), revealing persistent confusion between the two in the FEP literature. These critiques are highly relevant to EAS's core concern (how the boundary of $S \subset E$ is demarcated) but approach it from different angles.
\end{itemize}

\subsection{Divergence 4: Terminal State of Free Energy Minimization}

\begin{itemize}[leftmargin=*]
\item \textbf{FEP}: Non-equilibrium steady state (NESS); variational density approximates the true posterior.
\item \textbf{EAS}: Even upon reaching NESS, internal states still encode a quotient space of external states---$R_S$ rather than $E$ itself. Free energy minimization can only shrink the distance between quotient space and original space, not eliminate it.
\item \textbf{Friston's possible response}: FEP never claimed internal states fully encode external states---the ``as if'' framework makes no such commitment.
\item \textbf{EAS's follow-up}: If FEP does not commit to $q \to p$ convergence, then FEP's description of ``self-evidencing'' is not about truth but about survival---fundamentally different from EAS's epistemic concerns.
\end{itemize}

\subsection{Divergence 5: Status of Lean4 Formalization}

\begin{itemize}[leftmargin=*]
\item \textbf{EAS}: All theorems T0--T7 compile and pass in Lean4; formal verification ensures the \emph{inference chain} is unimpeachable (though the axiom itself can be questioned---see below).
\item \textbf{FEP}: FEP lacks equivalent formal verification. Biehl, Pollock \& Kanai (2021) proved by counterexample that FEP's Free Energy Lemma fails under certain conditions \cite{biehl2021}. \textbf{However}: Friston, Da Costa \& Parr (2021) published a formal response \cite{friston2021response}, arguing that Biehl et al.'s critique targets a ``particular formalization'' rather than FEP's core ideas. This debate remains unresolved.
\item \textbf{Friston's possible response}: FEP is a physical principle, not a mathematical theorem---analogous to the principle of least action; formal proof is not required.
\item \textbf{EAS's rebuttal}: Physical principles can be challenged (FEP has been revised multiple times), while Lean4-verified inference chains cannot be overturned. \textbf{However}, EAS's single axiom $S \subset E \implies R_S \not\cong E$ is itself not unquestionable---it is a metaphysical commitment about cognitive structure, not an empirically falsifiable proposition. Formalization ensures ``if one accepts the axiom, the consequences are inescapable,'' but does not ensure the axiom itself is true.
\end{itemize}

\section{Potential Dialogue Points}

\subsection{Common Language: Both Agree ``Cognition Is Bounded''}

FEP acknowledges exact inference is infeasible; EAS proves perfect self-modeling is impossible. Both share an intuition: \textbf{finite systems cannot fully grasp infinite reality}. The difference lies in the nature of this ``finitude'': FEP says ``computationally finite,'' EAS says ``structurally finite.''

\subsection{EAS Can Provide First-Principles Support for FEP}

Notably, Dietz's self-published paper (not peer-reviewed) \cite{dietz} also independently derives from philosophical first principles a conclusion structurally equivalent to FEP---``free energy is remainder; prediction error minimization is supersession.'' \textbf{Although this paper has not undergone peer review and its claim of ``structural equivalence'' remains to be tested by the scholarly community}, it represents a direction of parallel thinking. EAS's single axiom arrives at a similar position via a different path: \textbf{because $S \subset E \implies R_S \not\cong E$, finite systems must minimize free energy}---free energy minimization is the only survival strategy for quotient cognitive systems in NESS. The convergence of three frameworks (FEP, Dietz's closure framework, EAS) from different directions is itself worth further exploration.

\subsection{``Incompleteness'' as Completion of FEP, Not Negation}

EAS is not a replacement for FEP but its epistemic completion. FEP answers ``how cognitive systems operate''; EAS answers ``why cognitive systems cannot operate better.'' Together, they form a complete picture.

\section{Conclusion}

\begin{enumerate}[leftmargin=*]
\item \textbf{FEP has not yet theorem-ized the ``incompleteness'' problem.} The 2025 Fields--Friston paper touches on this topic (acknowledging systems cannot have fully precise models of their QRFs); Bruineberg et al.\ (2022) and Raja et al.\ (2021) critically examine the ontological status of Markov blankets from different angles; but none have elevated cognitive incompleteness to an insurmountable theorem. This is the core gap EAS can fill.

\item \textbf{FEP does not assume an ideal observer, but preserves the ideal observer as a mathematical limit.} EAS proves this limit is unreachable---not due to insufficient computational resources, but due to epistemic structure.

\item \textbf{Epiplexity five-layer resonance and variational free energy operate at different levels.} The former is a validation methodology; the latter is a descriptive principle. The former is more rigorous but harder to operationalize; the latter is more computable but lacks epistemic foundations.

\item \textbf{Optimal dialogue entry point}: Using EAS's single axiom to provide a first-principles proof for FEP's claim that ``approximate inference can only be approximate''---this is not challenging FEP, but completing its philosophical foundations.
\end{enumerate}

\newpage
\begin{thebibliography}{9}

\bibitem{friston2019}
Friston, K.J. (2019).
\textit{A free energy principle for a particular physics}.
arXiv:1906.10184.

\bibitem{friston2013}
Friston, K.J., Da Costa, L. \& Parr, T. (2013).
\textit{Life as we know it}.
Journal of the Royal Society Interface, 10(86), 20130475.

\bibitem{friston2025}
Friston, K.J., Ramstead, M.J.D. \& Sakthivadivel, A.T. (2025).
\textit{Generative modelling in non-equilibrium statistical mechanics}.
arXiv:2406.11630v3.

\bibitem{fields2025}
Fields, C., Albarracin, N., Friston, K.J. et al. (2025).
\textit{Inner screens and imaginative experience}.
Neuroscience of Consciousness, 2025(1), niaf009.
\url{https://academic.oup.com/nc/article/2025/1/niaf009/8117684}

\bibitem{laukkonen2025}
Laukkonen, R.E., Friston, K.J. \& Chandaria, S. (2025).
\textit{A Beautiful Loop}.
Neuroscience \& Bioreviews.
\url{https://doi.org/10.1016/j.neubiorev.2025.106296}

\bibitem{biehl2021}
Biehl, M., Pollock, F. \& Kanai, R. (2021).
\textit{A Technical Critique of Some Parts of the Free Energy Principle}.
Entropy, 23(3), 293.
\url{https://doi.org/10.3390/e23030293}

\bibitem{friston2021response}
Friston, K.J., Da Costa, L. \& Parr, T. (2021).
\textit{Some Interesting Observations on the Free Energy Principle}.
Entropy, 23(8), 1076.
\url{https://doi.org/10.3390/e23081076}

\bibitem{bruineberg2022}
Bruineberg, J., Dolega, K., Dewhurst, J. \& Baltieri, M. (2022).
\textit{The Emperor's New Markov Blankets}.
Behavioral and Brain Sciences, 45, e131.
\url{https://doi.org/10.1017/S0140525X21002351}

\bibitem{dietz}
Dietz, C.F. (n.d.).
\textit{The Grammar of Prediction: How the Closure Framework Grounds Karl Friston's Free Energy Principle}.
\url{https://carldietz.com/papers/downloads/Grammar_of_Prediction.pdf}

\end{thebibliography}

\end{document}
